\begin{result}{E.2.27}{A profinite loop with integral classes on every finite cover}{Proved}
Let \(X=\mathbf P^1(\C)\setminus\{0,1,\infty\}\), with basepoint
\(1/2\), and let \(F=\pi_1^{\mathrm{top}}(X,1/2)=\langle x,y\rangle\),
where \(x,y\) are the usual positive loops around \(0,1\).
Put \(G=\widehat F\), identified with the geometric \'{e}tale
fundamental group by the actual finite-cover comparison.
For an open subgroup \(H\subset G\), let \(Y_H\) be its pointed cover
and put
\[
 F_H=F\cap H,\qquad
 \Lambda_H=F_H^{\mathrm{ab}}=H_1(Y_H(\C),\Z),\qquad
 H^{\mathrm{ab}}=\Lambda_H\otimes\widehat\Z.
 \tag{E.2.27a}
\]
There is a fully specified element \(g\in G\) with the following
properties:
\begin{enumerate}
\item \(g\) is not conjugate in \(G\) to any element of \(F\).
\item For every open \(H\), every \(m\in\Z\), and every \(t\in G\)
such that \(tg^mt^{-1}\in H\),
\[
 \operatorname{ab}_H(tg^mt^{-1})\in\Lambda_H
       \subset\Lambda_H\otimes\widehat\Z.
 \tag{E.2.27b}
\]
\item \(g\neq1\), its image in \(G^{\mathrm{ab}}\) is zero, and
\(g\) lies in no conjugate of a cusp inertia subgroup.
\end{enumerate}
The image of (E.2.27b) under the actual filling map also belongs to
the integral homology lattice of the smooth compactification of
\(Y_H\). Thus these element-to-homology maps, even for all powers
and all finite-cover lifts, do not recognize a discrete conjugacy
class. The construction supplies an element of \(G\); it does not
supply the Galois automorphism and compatible full transporter
required in E.2.22.
\end{result}
\begin{proof}
We specify the finite choices, prove their compatibility with
commutator tails, and identify the exact integral representative
in each cover.

\paragraph{1. Fixed enumerations and the actual finite matrix maps.}
Order the letters \(x,x^{-1},y,y^{-1}\) as written, and enumerate
reduced words \(f_1,f_2,\ldots\) first by length and then
lexicographically, starting with \(f_1=1\). For each positive
integer \(a\), enumerate all ordered pairs of permutations of
\(\{1,\ldots,a\}\) lexicographically by their value tuples. Each
pair specifies a homomorphism
\(\phi_i:F\to S_{a_i}\), by the images of \(x,y\).
Retain every pair, including repetitions and nontransitive pairs,
and set
\[
 N_i=\ker(\phi_i:G\longrightarrow S_{a_i}).
 \tag{E.2.27c}
\]
Every finite quotient of \(F\) embeds in a symmetric group by its
regular action, so its kernel occurs in this list. Consequently
\(\bigcap_iN_i=1\), and for every open \(H\subset G\) some
\(N_i\) is contained in \(\operatorname{core}_G(H)\).

Define the homomorphism
\[
 \rho:F\longrightarrow\operatorname{SL}_2(\Z),\qquad
 \rho(x)=\begin{pmatrix}1&3\\0&1\end{pmatrix},\qquad
 \rho(y)=\begin{pmatrix}1&0\\3&1\end{pmatrix}.
 \tag{E.2.27d}
\]
Its faithfulness is not needed. Reduction modulo any positive
integer \(r\) extends to the actual finite homomorphism
\(\rho_r:G\to\operatorname{SL}_2(\Z/r\Z)\).
Every matrix in \(\rho(F)\) is \(I\) modulo \(3\), so its
integer trace is \(2\) modulo \(3\) and hence is nonzero.
With the convention \([u,v]=uvu^{-1}v^{-1}\), put
\(h_d=[x^d,y^d]\) for \(d\geq1\). Direct multiplication,
with \(z=3d\), gives
\[
 \rho(h_d)=
 \begin{pmatrix}1+z^2+z^4&-z^3\\z^3&1-z^2\end{pmatrix},
 \qquad \det\rho(h_d)=1,
 \qquad \operatorname{tr}\rho(h_d)=2+z^4>2.
 \tag{E.2.27e}
\]
Whenever \(x^d,y^d\) belong to an open normal subgroup \(W\),
the finite word \(h_d\) belongs to
\([F\cap W,F\cap W]\). This is the commutator freedom used
at every step below.

\paragraph{2. The induction and its three-term trace calculation.}
Set \(p_0=1\) and \(V_0=G\). At step \(n\) put
\(W_n=V_{n-1}\cap N_n\). Choose \(d_n\) to be the least positive
integer for which \(x^{d_n},y^{d_n}\in W_n\). Concretely,
\(d_n\) is the least common multiple of the orders of the two
generator images in \(\phi_i\) for \(i\leq n\), and in
\(\rho_{r_j}\) for \(j<n\), where the earlier moduli are defined
below. All these groups are finite, so this is a specified finite
integer. Write
\[
 h_n=h_{d_n},\qquad P=\rho(p_{n-1}),\qquad B=\rho(h_n),
 \qquad c=\operatorname{tr}\rho(f_n),\qquad T=\operatorname{tr}B.
 \tag{E.2.27f}
\]
At least one of \(m=0,1,2\) satisfies
\(\operatorname{tr}(PB^m)\neq c\). Indeed, if all three traces
equalled \(c\), the Cayley--Hamilton identity
\(B^2-TB+I=0\), multiplied by \(P\) and traced, would give
\[
 c-Tc+c=0,\qquad (T-2)c=0.
 \tag{E.2.27g}
\]
But \(T>2\) by (E.2.27e), while the equality at \(m=0\) would
give \(c=\operatorname{tr}P\equiv2\pmod3\), so \(c\neq0\).
This proves that the step cannot stall. Choose the least such
\(m=m_n\), and set
\[
 \begin{split}
 p_n&=p_{n-1}h_n^{m_n},\\
 \Delta_n&=\operatorname{tr}\rho(p_n)-\operatorname{tr}\rho(f_n)
                    \neq0,\qquad r_n=|\Delta_n|+1,\\
 V_n&=W_n\cap\ker\rho_{r_n}.
 \end{split}
 \tag{E.2.27h}
\]
The nonzero integer \(\Delta_n\) is divisible by \(3\), so
\(r_n\geq4\), and its residue modulo \(r_n\) is nonzero.
Composite moduli cause no problem: trace is defined and invariant
under conjugation over \(\Z/r_n\Z\).

Since \(h_n\in W_n\subset V_{n-1}\), the nonempty compact cosets
\(p_nV_n\) are nested. Their intersection is a singleton:
\(V_n\subset\bigcap_{i\leq n}N_i\) and
\(\bigcap_iN_i=1\). Define \(g\) to be this unique element.
Equivalently, the finite words \(p_n\) converge to \(g\) in
\(G\). All enumerations, least choices, and moduli in this
definition have been fixed explicitly.

\paragraph{3. Exclusion of every entire profinite conjugacy class.}
For every \(n\), membership \(g\in p_nV_n\) implies
\[
 \operatorname{tr}\rho_{r_n}(g)
   =\operatorname{tr}\rho_{r_n}(p_n)
   \neq\operatorname{tr}\rho_{r_n}(f_n).
 \tag{E.2.27i}
\]
Conjugation in \(G\) induces conjugation under this actual finite
matrix map, which preserves trace. Thus \(g\) is not conjugate
to \(f_n\). Every discrete word occurs, including \(1\), so
property 1 and \(g\neq1\) follow. Each exclusion is preserved by
all later tails because those tails lie in the frozen kernel
\(\ker\rho_{r_n}\). No conjugacy-separation theorem is required.

\paragraph{4. The prefix that lies in each finite cover.}
Fix an open \(H\subset G\), and choose \(i\) with
\(N_i\subset\operatorname{core}_G(H)\). Put
\[
 K_i=F\cap N_i,\qquad C_i=\overline{[K_i,K_i]}^{\,G}.
 \tag{E.2.27j}
\]
The subgroup \(K_i\) is normal in \(F\), as is its commutator
subgroup. Density of \(F\) and continuity of conjugation show
that \(C_i\) is normal in \(G\). Moreover,
\(K_i\subset F_H\), so
\[
 C_i\subset\ker(H\longrightarrow H^{\mathrm{ab}})\subset H.
 \tag{E.2.27k}
\]
For every \(j\geq i\), the inclusion \(W_j\subset N_i\) gives
\(h_j^{m_j}\in[K_i,K_i]\). Taking the limit of the actual tails
therefore gives
\[
 g=p_{i-1}c_i\quad(c_i\in C_i),\qquad
 g^mC_i=p_{i-1}^{,m}C_i\quad(m\in\Z).
 \tag{E.2.27l}
\]
The second equality holds in the group \(G/C_i\), using the
normality just proved; it includes negative powers. Whenever
\(g^m\in H\), it follows that \(p_{i-1}^{,m}\in H\cap F\).
Their images in \(H^{\mathrm{ab}}\) are equal. This argument
does not require \(p_{i-1}\) itself to belong to \(H\).

Now let \(t\in G\) and suppose \(tg^mt^{-1}\in H\). The open
left coset \(Ht\) meets \(F\); choose \(a\in F\cap Ht\).
Then \(t=ha\) for some \(h\in H\), and conjugation by \(h\)
acts trivially on \(H^{\mathrm{ab}}\). Normality of \(C_i\)
and (E.2.27l) give
\[
 ag^ma^{-1}C_i=ap_{i-1}^{,m}a^{-1}C_i.
 \tag{E.2.27m}
\]
The left representative is in \(H\); since \(C_i\subset H\),
the right representative lies in \(H\cap F=F_H\).
Thus the actual integral representative is explicitly identified:
\[
 \operatorname{ab}_H(tg^mt^{-1})
   =\operatorname{ab}_H(ap_{i-1}^{,m}a^{-1})\in\Lambda_H.
 \tag{E.2.27n}
\]
This proves (E.2.27b) for every profinite conjugate lift.

For completeness, the finite-index subgroup \(F_H\subset F\)
has its full profinite topology: the core in \(F\) of any
finite-index subgroup of \(F_H\) is finite-index and contained
in that subgroup. Hence its completion is \(H\). The finite
abelian quotient universal property then identifies
\(H^{\mathrm{ab}}\) with the completion of \(F_H^{\mathrm{ab}}\).
The latter is a finite-rank free abelian group, giving exactly
(E.2.27a). The topological/
\'{e}tale finite-cover comparison is the one used in E.2.22
\cite[Corollary 3.7, p.~10]{OSU2024}.
The continuous map from \(Y_H(\C)\) to its compactification
induces the integral homology map and its completion. Applying
it to (E.2.27n) proves the stated proper-homology assertion.

\paragraph{5. Cusp inertia and the remaining transporter requirement.}
Every \(p_n\) is a product of commutators, so the image of
\(g\) in \(G^{\mathrm{ab}}=\widehat\Z^{,2}\) is zero.
The three inertia generators \(x,y,(xy)^{-1}\) have primitive
homology vectors. If \(g\) were conjugate to such a generator
raised to \(\alpha\in\widehat\Z\), its homology would be
\(\alpha\) times that primitive vector. This forces
\(\alpha=0\), hence \(g=1\), which (E.2.27i) excludes.
Property 3 follows. By F.3.33, the nontrivial procyclic subgroup
\(\overline{\langle g\rangle}\) therefore has no arithmetic
normalizer with open Galois image.

To compare the basepoint with E.2.22, use the upper semicircle
with centre \(5/4\) and radius \(3/4\), oriented from \(2\) to
\(1/2\). This path avoids \(0,1\) and sends a loop \(\gamma\)
at \(1/2\) to \(\beta\gamma\beta^{-1}\) at \(2\).
The induced discrete and profinite isomorphisms carry the
corresponding covers and their integral homology maps to each
other. Thus the element-to-homology limitation applies to the
actual tower in E.2.22. That result additionally prescribes one
Galois automorphism on the entire tower and all compatible maps;
the present construction supplies neither that automorphism nor
a full motivic transporter, and leaves that exact requirement open.

The standalone certificate
\path{code/check_nondiscrete_integral_loop.py} verifies the
commutator matrix, determinant, trace, Cayley--Hamilton identity,
and the first \(96\) specified induction stages using exact
integers and finite permutations. It checks preservation of all
earlier matrix cylinders and trace exclusions. The compact limit
and all-cover claims are proved above and are not inferred from
those finite checks.
\end{proof}
